\documentclass[11pt]{article}
\usepackage[margin=1in]{geometry}
\usepackage{amsmath,amssymb,amsthm}
\usepackage{enumitem}

\title{Proposed improvements to \texttt{EpidCRN} based on SubFF.pdf}
\author{Generated from SubFF.pdf analysis}
\date{February 2026}

\begin{document}
\maketitle

\section{Recipe~II implementation (\texttt{oscRII})}

The paper defines (Definition~4, item~2): \emph{An Oscillatory Core of Class~II
is a minimal $k \times k$ CS-matrix which is an unstable negative feedback, and
which has a $(k-1) \times (k-1)$ Hurwitz-stable principal submatrix.}

Currently only Recipe~I (\texttt{oscRI}) is implemented. The analogous function
\texttt{oscRII} would:
\begin{enumerate}[nosep]
  \item Select all UNF nodes (using \texttt{unfQ}).
  \item For each UNF of size $k$, enumerate its $(k-1)\times(k-1)$ principal
        submatrices (delete one row and the corresponding column).
  \item Check Hurwitz stability of each submatrix via \texttt{hurQ}.
  \item Keep minimal UNFs (no proper UNF sub-CS also satisfies the condition).
\end{enumerate}

Proposed Mathematica code (for \texttt{Core.wl}, after \texttt{oscRI}):
\begin{verbatim}
(* oscRII: Recipe II -- Oscillatory Cores of Class II.
   A core is a UNF (unstable negative feedback) of size k
   that has a (k-1)x(k-1) Hurwitz-stable principal submatrix.
   Keep only minimal such UNFs. *)
oscRII[nCS_, gH_, spe_, rts_] := Module[
  {unfIdx, cores, k, Sk, subMats},
  unfIdx = Select[Range[Length[nCS]], unfQ[nCS[[#]]] &];
  cores = Select[unfIdx, Function[u,
    k = nCS[[u, 1]]; Sk = nCS[[u, 5]];
    If[k < 2, False,
      subMats = Table[
        Sk[[Delete[Range[k], j], Delete[Range[k], j]]],
        {j, k}];
      AnyTrue[subMats, hurQ]]]];
  (* minimality: remove if a smaller UNF core is a sub-CS *)
  cores = Select[cores, Function[d,
    !AnyTrue[Complement[cores, {d}],
      MemberQ[VertexOutComponent[gH, #], d] &]]];
  If[Length[cores] == 0,
    Print["Recipe II: no oscillatory cores of Class II found"],
    Print["Recipe II: ", Length[cores], " cores of Class II"];
    Do[Module[{c = cores[[j]]},
      Print["RII#", j,
        " species=", spe[[nCS[[c, 3]]]],
        " reactions=", rts[[nCS[[c, 4]]]],
        " Sk=", nCS[[c, 7]] // MatrixForm]],
      {j, Length[cores]}]];
  cores];
\end{verbatim}


\section{Recipe~0: fixed-sign determinant test (\texttt{fixSgn})}

Corollary~5.7/5.8 of Vassena states: the coefficient $c_k$ is of fixed sign
for all positive $\mathbf{r}'$ iff all $k$-CS have sign-modified determinants of the
same sign. If the Jacobian $G$ has a nonzero determinant of fixed sign, and the
network admits both stability and instability, then $G(\mathbf{r}')$ has purely
imaginary eigenvalues (Recipe~0).

Currently checked manually in \texttt{SIRWS.wl} (lines 84--89). Proposed
function:
\begin{verbatim}
(* fixSgn: test fixed-sign condition for each coefficient c_k.
   Returns list of {k, isFixedSign, allSigns} for k=1..n.
   If c_n has fixed sign AND UPFs exist AND stable n-CS exist -> Hopf. *)
fixSgn[nCS_, n_] := Module[{csK, dets, sgns},
  Table[
    csK = Select[nCS, #[[1]] == k &];
    dets = #[[6]] & /@ csK;
    sgns = Sign /@ dets;
    {k, And @@ (# >= 0 & /@ dets) || And @@ (# <= 0 & /@ dets),
     Union[sgns]},
    {k, n}]];
\end{verbatim}

Proposed script addition (for the notebooks, after \texttt{negCS}):
\begin{verbatim}
fs = fixSgn[nCS, n];
Print["Fixed sign by order: ",
  Table[{k, fs[[k, 2]]}, {k, n}]]
\end{verbatim}


\section{Symbolic Hopf hunt via resultants (\texttt{hopfRes})}

Section~5.4 describes the ``symbolic hunt'' method for Hopf bifurcations:
\begin{enumerate}[nosep]
  \item Compute characteristic polynomial $P(\lambda) = \det(G - \lambda\,\mathrm{Id})$.
  \item Evaluate at $\lambda = i\omega$ and split into real and imaginary parts.
  \item Divide imaginary part by $\omega$ (removing trivial root $\omega=0$).
  \item Set $\Omega = \omega^2$ and compute $\mathrm{resultant}(\Re(P),\,\tilde\Im(P),\,\Omega)$.
  \item Hopf occurs where the resultant vanishes with $\Omega > 0$.
\end{enumerate}

Proposed Mathematica function:
\begin{verbatim}
(* hopfRes: symbolic Hopf hunt via resultant of Re/Im parts.
   G = symbolic Jacobian (n x n matrix, entries may contain symbols).
   Returns {resultant, Omega=omega^2, ReP, ImPtilde}. *)
hopfRes[G_] := Module[{n, cp, PG, rP, iP, iPt, Om, res},
  n = Length[G];
  cp = CharacteristicPolynomial[G, lam];
  PG = cp /. lam -> I om;
  rP = ComplexExpand[Re[PG]];
  iP = ComplexExpand[Im[PG]];
  iPt = Cancel[iP/om];
  Om = om^2;
  rP = rP /. om^k_?EvenQ :> Om^(k/2);
  iPt = iPt /. om^k_?EvenQ :> Om^(k/2);
  res = Resultant[rP, iPt, Om];
  {Simplify[res], Om, rP, iPt}];
\end{verbatim}

This directly implements equations (5.29)--(5.30) of the paper.


\section{Necessary conditions for bifurcation (\texttt{bifNec})}

Theorem~4 and Lemma~3 prove that for any bifurcation at an endemic
equilibrium $(\bar s,\bar\imath,\bar r) > 0$, it is necessary that
$G_{ii}(\bar s,\bar\imath) > 0$, equivalently
$f'_i(\bar s,\bar\imath) > h'(\bar\imath) + \mu_i$.

Observation~3: if $f'_i \le 0$ then $G$ is weakly-diagonally dominant and no
bifurcation is possible. This excludes decreasing infection rates.

These are analytic results, but could be checked numerically at a given
equilibrium:
\begin{verbatim}
(* bifNec: check necessary condition G_ii > 0 at endemic equilibrium.
   jac = Jacobian matrix, iIdx = index of infected variable.
   Returns True if necessary condition is satisfied. *)
bifNec[jac_, iIdx_] := jac[[iIdx, iIdx]] > 0;
\end{verbatim}


\section{Purely imaginary eigenvalue theorem for ME (\texttt{meHopf})}

Theorem~2 proves that any epidemiological system including reactions
$S+I \to 2I$ (infection) and $I \to \ldots$ (recovery/death) admits purely
imaginary eigenvalues of the symbolic Jacobian.

This is a structural test that could be automated: given the reaction network
representation, check if it contains an autocatalytic infection reaction and a
removal reaction for $I$:
\begin{verbatim}
(* meHopf: check if ME model structurally admits Hopf.
   Requires: infection S+I -> 2I (autocatalysis) present in RN,
   and a removal reaction I -> ... (not producing S or I).
   By Theorem 2, symbolic Jacobian then admits purely imaginary eigenvalues. *)
meHopf[RN_, var_, alp_, bet_] := Module[{n, sIdx, iIdx, autoQ, remQ, j},
  n = Length[var];
  (* find S and I candidates: I is autocatalytic (alp and bet both > 0),
     S is consumed (alp > 0, bet = 0 or bet < alp) *)
  autoQ = False; remQ = False;
  Do[
    If[alp[[iIdx, j]] > 0 && bet[[iIdx, j]] > alp[[iIdx, j]],
      (* autocatalytic in species iIdx for reaction j *)
      autoQ = True],
    {iIdx, n}, {j, Length[RN]}];
  Do[
    If[alp[[iIdx, j]] > 0 && bet[[iIdx, j]] == 0 &&
       Total[bet[[All, j]]] > 0,
      remQ = True],
    {iIdx, n}, {j, Length[RN]}];
  autoQ && remQ];
\end{verbatim}


\section{Equilibrium flux vector (\texttt{eqFlux})}

The paper (equation~5.6) emphasizes that the existence of a positive
equilibrium requires a positive right kernel vector $\bar{\mathbf{v}} > 0$
of $\Gamma$, called an equilibrium flux vector. This is a prerequisite
check before attempting bifurcation analysis:
\begin{verbatim}
(* eqFlux: find equilibrium flux vector v > 0 such that gam.v = 0.
   Returns {} if no positive kernel vector exists (deficiency 0),
   or a symbolic parametrization of the kernel. *)
eqFlux[gam_] := Module[{ker},
  ker = NullSpace[gam];
  If[ker === {}, Print["No kernel: rank = dimension"]; Return[{}]];
  ker];
\end{verbatim}

This is useful for the scripts: if $\Gamma$ has full column rank,
no positive equilibrium can exist in the stoichiometric compatibility class,
which rules out endemic equilibria.


\section{Commutative diagram: projection to Boros--Rost witness (\texttt{projBR})}

Section~4.8 describes the commutative diagram:
\[
\mathcal{N}_{\mathrm{SIRWS}} \xrightarrow{\pi} \mathcal{N}_{\mathrm{BR}}
\]
where $\pi$ projects out intermediaries (e.g.\ $R$) and collapses pathways.
Given a Recipe~I core $(\kappa, \kappa')$, the projection constructs a
reduced mass-action network (the Boros--Rost witness).

This is model-specific and difficult to automate in general, but the
\emph{detection} of the core structure is already done by \texttt{oscRI}.
A useful addition would be to print the reduced stoichiometric matrix
of the core alongside the full system, to help identify the projection
manually:
\begin{verbatim}
(* In oscRI, after printing each core, also print the
   block structure of the full n x n CS containing it: *)
(* Already partially done: oscRI prints Sk for both core and UPFs.
   Improvement: print eigenvalues of the core Sk block to show
   the oscillatory pair explicitly. *)
Do[Module[{c = cores[[j]], eigC},
  eigC = Eigenvalues[N[nCS[[c, 5]]]];
  Print["RI#", j, " eigenvalues=", eigC]],
  {j, Length[cores]}];
\end{verbatim}


\section{Improve \texttt{upfCS} output with determinant and size}

Currently \texttt{upfCS} prints species and rates but not the determinant sign
or the CS order $k$. The paper emphasizes that the sign
$\sigma(A) = \mathrm{sign}((-1)^k \det A)$ determines the type of instability
(Lemma~1). Adding the determinant and order $k$ to the print would help match
the paper's analysis:
\begin{verbatim}
(* In upfCS, change the print to include k and det: *)
Print["UPF#", j,
  " k=", nCS[[ui, 1]],
  " species=", spe[[nCS[[ui, 3]]]],
  " rates=", rts[[nCS[[ui, 4]]]],
  " det=", nCS[[ui, 6]],
  " stable supersets=", spe[[nCS[[#, 3]]]] & /@ di]
\end{verbatim}


\section{Notation alignment with paper}

The paper uses specific notation that differs from the code in some places.
Aligning would help readers cross-reference:

\begin{center}
\begin{tabular}{lll}
\textbf{Paper} & \textbf{Current code} & \textbf{Proposed change} \\
\hline
$\Gamma$ (stoichiometric matrix) & \texttt{gam} & keep \\
$\alpha$ (reactant matrix) & \texttt{alp} & keep \\
$\mathcal{R}$ (reactivity matrix) & \texttt{RM} & keep \\
$G = \Gamma \cdot \mathcal{R}$ (symbolic Jacobian) & \texttt{G} & keep (in scripts) \\
$c_k$ (char poly coefficients) & \texttt{co[[k]]} & keep \\
$\alpha_{\mathbf{J}^{(k)}} = \det S[\mathbf{J}^{(k)}]$ & \texttt{nCS[[i,6]]} & add comment: ``sign-modified det'' \\
$\sigma(A) = \mathrm{sign}((-1)^k \det A)$ & \texttt{upfQ}/\texttt{unfQ} & add comment referencing Lemma~1 \\
$(\kappa, E_\kappa, J)$ (child selection) & \texttt{csInfo} output & document: \texttt{sIdx}$=\kappa$, \texttt{rBij}$=J$ \\
\end{tabular}
\end{center}

Concrete: update the \texttt{csInfo} comment header to read:
\begin{verbatim}
(* csInfo: given order k, a negative term from co[[k]], and gam,
   extract the k-child selection (kappa, E_kappa, J):
   - kappa = sIdx (species subset),
   - J = rBij (bijection: species -> reactions, in species order),
   - S[J] = Sk = gam[[sIdx, rIdx]] (stoich submatrix, sorted columns),
   - alpha_J = Det[Sk] (sign-modified determinant),
   - SkB = gam[[sIdx, rBij]] (bijection-ordered, for display).
   Returns {k, term, sIdx, rBij, Sk, Det[Sk], SkB}. *)
\end{verbatim}


\section{Add \texttt{oscRII} usage to \texttt{EpidCRN.wl}}

When \texttt{oscRII} is added to \texttt{Core.wl}, its usage statement
must be added to \texttt{EpidCRN.wl}:
\begin{verbatim}
oscRII::usage = "oscRII[nCS,gH,spe,rts]: Recipe II oscillatory cores";
\end{verbatim}

Similarly for any new functions (\texttt{fixSgn}, \texttt{hopfRes}, etc.).


\section{Summary of priorities}

\begin{enumerate}
  \item \textbf{Recipe~II} (\texttt{oscRII}): highest priority, directly extends
    the existing Recipe~I and completes the oscillatory core classification from
    Definition~4.
  \item \textbf{Fixed-sign test} (\texttt{fixSgn}): simple, replaces hardcoded
    checks in notebooks.
  \item \textbf{Hopf resultant} (\texttt{hopfRes}): implements the symbolic hunt
    from Section~5.4, useful for the SIRScrw model analysis.
  \item \textbf{Notation/comment updates}: low effort, high clarity.
  \item \textbf{\texttt{upfCS} output enhancement}: minor, adds $k$ and $\det$.
  \item \textbf{ME Hopf structure test} (\texttt{meHopf}): useful but narrow
    (only confirms what Theorem~2 guarantees for ME models).
  \item \textbf{Equilibrium flux vector}: informational, already implicitly
    checked in NGMRN.
  \item \textbf{Bifurcation necessary condition}: analytic check, useful for
    the SIRScrw analysis.
\end{enumerate}

\end{document}
